% Lösungen Einheit 1 von 4 – Rechte Winkel an Dreiecken nachweisen (zusammenbau v0.1)
\einheitenkopf{Einheit 1 von 4 · Rechte Winkel an Dreiecken nachweisen}

\erg{9}{a) Gerade gegen Ebene – Richtungsvektor kollinear zum Normalenvektor}
\erg{10}{a) null zeigen – einsetzen und ausrechnen}
\erg{11}{a) für alle – das Skalarprodukt muss für jeden Wert null sein}
\erg{12}{a) $\overrightarrow{BA}$ und $\overrightarrow{BC}$ – beide Schenkel gehen vom Scheitel $B$ aus \quad b) $\overrightarrow{AB} = (3 | 2 | 0)$, $\overrightarrow{AC} = (-2 | 3 | 0)$; $\overrightarrow{AB} \circ \overrightarrow{AC} = 3 \cdot (-2) + 2 \cdot 3 + 0 \cdot 0 = 0$ – rechter Winkel bei $A$ \quad c) $\overrightarrow{OA} \circ \overrightarrow{OB} = 0 + 0 + 0 = 0$; $|\overrightarrow{OB}| = 13$, $|\overrightarrow{OA}| = \sqrt{25 + 144} = 13$ \quad d) $\overrightarrow{AB} \circ \overrightarrow{AC_t} = 4 \cdot 3t + 3 \cdot (-4t) + 5 \cdot 0 = 12t - 12t = 0$ für jedes $t$ \quad e) $AB$ liegt in der $xy$-Ebene, $AC$ auf der $z$-Achse, die auf dieser Ebene senkrecht steht – rechter Winkel bei $A$ \quad f) bei $D$: $(0 | 20 | 0) \circ (-18 | 4 | 0) = 80$; bei $E$: $(-18 | -16 | 0) \circ (0 | -20 | 0) = 320$; bei $F$: $(18 | -4 | 0) \circ (18 | 16 | 0) = 260$ – keines der drei Skalarprodukte ist null, also gibt es keinen rechten Winkel}
\erg{13}{a) Fehler: Tom prüft die Ecke $A$ statt des behaupteten Scheitels $S$. Richtig: $\overrightarrow{SA} \circ \overrightarrow{SB} = (1 | 2 | 2) \circ (2 | -1 | 0) = 2 - 2 + 0 = 0$ – rechter Winkel bei $S$ \quad b) Der Winkel bei $B$ wird von den beiden Seiten gebildet, die in $B$ zusammenlaufen; ihre Richtungen sind die Vektoren von $B$ zu den anderen Ecken. $\overrightarrow{AB}$ und $\overrightarrow{AC}$ beschreiben den Winkel bei $A$. \quad c) $\overrightarrow{BA} \circ \overrightarrow{BC} = (-12) \cdot 0 + (-5) \cdot 0 + 0 \cdot 9 = 0$ – rechter Winkel bei $B$; $|\overrightarrow{AC}| = \sqrt{144 + 25 + 81} = \sqrt{250} \approx 15{,}8$ dm}
